\chapter{Discectomy}

\section*{Introduction \& Clinical Indications}
A discectomy is a surgical procedure primarily aimed at removing a portion of an intervertebral disc that is compressing a nerve root or the spinal cord. This intervention is most commonly performed in the lumbar spine, but can also be indicated for cervical or thoracic disc herniations. The surgery directly addresses symptoms such as radicular pain, numbness, weakness, or myelopathy, which arise from the mechanical compression and chemical irritation caused by the herniated disc material. It is a highly effective treatment for carefully selected patients who have failed conservative management or present with acute neurological deficits, significantly improving quality of life by alleviating pain and restoring neurological function. The goal is to decompress neural structures and facilitate their recovery.

\begin{itemize}
  \item Microdiscectomy
  \item Endoscopic Discectomy
  \item Open Discectomy
  \item Disc Excision
  \item Nerve Root Decompression
\end{itemize}

The fundamental principle of a discectomy is the direct removal of the offending disc material that is impinging upon neural structures. A herniated intervertebral disc, typically the nucleus pulposus, can protrude or extrude into the spinal canal or neural foramen, causing mechanical compression of the spinal cord or exiting nerve roots. This compression is often compounded by an inflammatory response from chemical irritants released by the disc, leading to symptoms such as radicular pain, paresthesia, and motor weakness.

The surgical rationale is to achieve immediate and sustained decompression of the neural elements. This is accomplished by carefully accessing the spinal canal, identifying the herniated disc fragment, and excising it using specialized instruments. The technical goals include complete removal of the symptomatic fragment while preserving the integrity of the surrounding healthy disc, minimizing trauma to the nerve root, and maintaining spinal stability. By relieving both the mechanical pressure and the chemical irritation, the procedure aims to resolve the patient's neurological symptoms and facilitate nerve recovery.

The surgical treatment of disc herniation has evolved significantly since its initial descriptions. The seminal work by Mixter and Barr in 1934 is widely credited with establishing the link between lumbar disc herniation and sciatica, and describing the first successful surgical removal of a herniated disc. Their technique involved a wide laminectomy to access the disc, which was associated with significant muscle dissection and postoperative pain.

Throughout the mid-20th century, open laminectomy and discectomy remained the standard. A major advancement occurred in the 1970s with the popularization of microsurgical discectomy by surgeons like M. Gazi Yasargil and Wolf-Dietrich Caspar. The introduction of the operating microscope allowed for smaller incisions, enhanced visualization of delicate neural structures, and reduced tissue dissection, leading to less postoperative pain and faster recovery. Subsequent innovations have included endoscopic and percutaneous techniques, further minimizing invasiveness and expanding the surgical options for disc herniation, aiming for even smaller incisions and quicker rehabilitation.

A discectomy is typically indicated for specific clinical scenarios where conservative management has failed or when there is an urgent neurological compromise.

\begin{itemize}
  \item Persistent radicular pain (sciatica in the leg or brachialgia in the arm) that has not improved after 6-12 weeks of appropriate non-surgical treatment, including physical therapy, medications, and epidural steroid injections.
  \item Progressive neurological deficit, such as worsening motor weakness (e.g., foot drop, wrist drop) or significant sensory loss, indicating ongoing nerve damage and potential for irreversible injury.
  \item Cauda equina syndrome, characterized by bilateral radiculopathy, saddle anesthesia, bowel or bladder dysfunction, and severe motor weakness, which constitutes a surgical emergency requiring urgent decompression.
  \item Myelopathy due to cervical or thoracic disc herniation, presenting with gait disturbance, spasticity, or upper motor neuron signs, indicating spinal cord compression.
  \item Severe, intractable pain that significantly impairs the patient's quality of life and functional capacity, despite maximal conservative efforts and clear correlation with imaging.
  \item Imaging evidence (typically MRI) of a disc herniation that clearly correlates with the patient's clinical symptoms and neurological findings, confirming the anatomical cause of the symptoms.
  \item Recurrent disc herniation at a previously treated level, if symptoms are severe and correlate with new imaging findings of neural compression.
\end{itemize}

Discectomy can serve as both a primary and a secondary surgical intervention, depending on the clinical context. It is considered a primary surgical treatment in cases of acute, severe neurological deficits, such as cauda equina syndrome or rapidly progressive motor weakness, where immediate decompression is critical to prevent permanent nerve damage. For patients experiencing severe, debilitating radicular pain that has failed to respond to a prolonged course of conservative management (typically 6-12 weeks), discectomy is also considered a primary surgical option, as it directly addresses the mechanical cause of their symptoms and offers a high likelihood of rapid pain relief.

Conversely, discectomy may be a secondary procedure in situations where a patient has previously undergone non-surgical interventions, such as epidural steroid injections or extensive physical therapy, without adequate relief. It can also be a secondary or salvage procedure for recurrent disc herniation at the same level after a prior discectomy, provided there is clear evidence of new compression and correlating symptoms. In general, it is not the first-line treatment for mild or resolving symptoms, as many disc herniations can improve with conservative care over time.

\section*{Relevant Surgical Anatomy}
The spine is a complex structure composed of vertebral bodies separated by intervertebral discs, which act as shock absorbers. Each disc consists of an outer fibrous ring, the annulus fibrosus, and a gel-like inner core, the nucleus pulposus. Posterior to the vertebral bodies and discs lies the spinal canal, which houses the spinal cord (terminating as the conus medullaris around L1-L2 in adults) and the cauda equina (a bundle of nerve roots below the conus). Nerve roots exit the spinal canal through the neural foramen, formed by adjacent pedicles and vertebral bodies. Key ligaments include the posterior longitudinal ligament (PLL), which runs along the posterior aspect of the vertebral bodies, and the ligamentum flavum, which connects adjacent laminae. A herniated disc typically protrudes posterolaterally, compressing the exiting nerve root or, less commonly, the spinal cord itself.

In the lumbar spine, the most common site for discectomy, the nerve root exiting at a given level (e.g., L5 nerve root at L4-L5 disc space) is typically compressed. Surgical landmarks include the spinous processes, laminae, and facet joints. The arterial supply to the spinal cord and nerve roots is via segmental arteries, which give rise to radicular arteries. Venous drainage occurs through an extensive epidural venous plexus. In the cervical spine, an anterior approach often involves navigating around the carotid sheath laterally and the esophagus/trachea medially. The uncinate processes are bony projections that help stabilize the cervical discs, and the vertebral arteries pass through the transverse foramina of the cervical vertebrae, requiring careful attention during anterior cervical discectomy.

\section*{Preoperative Workup \& Preparation}
Thorough preoperative preparation is crucial to ensure patient safety and optimize surgical outcomes for discectomy.

A comprehensive evaluation begins with a detailed review of the patient's medical history, current medications, allergies, and a thorough neurological examination to document baseline deficits. Imaging studies are paramount, with Magnetic Resonance Imaging (MRI) of the affected spinal region being the gold standard for visualizing disc herniations and nerve compression. X-rays may be obtained to assess spinal alignment and stability, and Computed Tomography (CT) scans can provide additional detail on bony anatomy, especially if fusion is contemplated.

\begin{itemize}
  \item \textbf{Laboratory Tests:} Routine preoperative blood tests typically include a complete blood count (CBC), basic metabolic panel (BMP), coagulation profile (PT/INR, PTT), and often a type and screen in case of significant blood loss.
  \item \textbf{Cardiac and Anesthesia Clearance:} Patients with significant comorbidities, especially cardiac or pulmonary conditions, will require clearance from their primary care physician or a specialist, along with a preoperative anesthesia evaluation to assess fitness for general anesthesia.
  \item \textbf{Medication Adjustments:} Patients on anticoagulants (e.g., warfarin, novel oral anticoagulants) or antiplatelet agents (e.g., aspirin, clopidogrel) will need to discontinue these medications for a specified period before surgery, as directed by the surgeon and anesthesiologist, to minimize bleeding risk.
  \item \textbf{NPO Status:} Patients are instructed to be NPO (nil per os - nothing by mouth) for a specific duration (typically 6-8 hours) before surgery to prevent aspiration during anesthesia.
  \item \textbf{Preoperative Antibiotics:} Prophylactic intravenous antibiotics are administered shortly before the incision to reduce the risk of surgical site infection.
  \item \textbf{Informed Consent:} A detailed discussion with the surgeon regarding the procedure, potential benefits, risks, alternatives, and expected recovery is conducted, and the patient provides informed consent, understanding the implications of the surgery.
\end{itemize}

While discectomy is a common and effective procedure, certain conditions may contraindicate its performance or necessitate specific precautions.

\textbf{Situations requiring diagnostic clarification, optimization, or an alternative plan:}
\begin{itemize}
  \item Active systemic infection or sepsis, which significantly increases the risk of surgical site infection, meningitis, and systemic complications.
  \item Uncontrolled coagulopathy or severe bleeding disorder that cannot be corrected preoperatively, posing an unacceptable risk of hemorrhage.
  \item Severe medical comorbidities (e.g., unstable angina, recent myocardial infarction, severe respiratory failure) that preclude safe administration of general anesthesia.
  \item Local skin infection at the planned surgical site, which could contaminate the surgical field.
  \item Lack of clear correlation between imaging findings and clinical symptoms, suggesting the disc herniation may not be the cause of the patient's pain.
\end{itemize}

\textbf{Relative Contraindications \& Precautions:}
\begin{itemize}
  \item Mild, non-progressive symptoms that are still responding to conservative management, as many disc herniations resolve spontaneously.
  \item Significant psychological overlay or somatization disorder, which may impact perceived outcomes and patient satisfaction despite technically successful surgery.
  \item Severe spinal instability requiring a fusion procedure, as discectomy alone may not be sufficient to address the underlying biomechanical issue.
  \item Previous failed surgery at the same level without clear evidence of new pathology or recurrent compression, necessitating careful re-evaluation.
  \item Obesity, which can increase surgical complexity, prolong operative time, and elevate the risk of complications such as wound infection and deep vein thrombosis.
  \item Smoking, which is associated with higher rates of complications, poorer wound healing, and increased risk of recurrent disc herniation.
\end{itemize}
\textbf{Precautions:} Intraoperative verification of the correct spinal level using fluoroscopy or other imaging modalities is paramount to prevent wrong-level surgery. Careful handling of neural structures, meticulous hemostasis, and avoidance of excessive retraction are also critical to minimize complications.

\section*{Surgical Technique \& Procedure}
Discectomy is performed under general anesthesia, with the patient carefully positioned to optimize surgical access and minimize complications. For lumbar discectomy, the patient is typically placed in a prone position on a specialized spinal frame to flex the lumbar spine, which helps open the interlaminar space and reduce epidural venous bleeding. For cervical discectomy, an anterior approach usually involves a supine position with the neck extended and slightly rotated.

The procedure generally follows these key steps:
\begin{itemize}
  \item \textbf{Anesthesia and Positioning:} General anesthesia is administered. The patient is meticulously positioned, and the surgical site is prepped and draped in a sterile fashion. Intraoperative fluoroscopy or navigation systems are routinely used to confirm the correct spinal level before incision.
  \item \textbf{Incision and Exposure:} For a lumbar microdiscectomy, a small midline incision (typically 2-4 cm) is made over the affected spinal level. For an anterior cervical discectomy, a small transverse incision is made in a skin crease on the front of the neck. The muscles are carefully dissected and retracted to expose the lamina and facet joint complex (lumbar) or the anterior vertebral bodies and disc space (cervical).
  \item \textbf{Bone Removal (Laminotomy/Foraminotomy):} A small window of bone (laminotomy or partial hemilaminectomy) may be removed from the lamina and/or medial facet to gain access to the spinal canal and nerve root, minimizing disruption to spinal stability. For cervical discectomy, the entire disc is removed, often followed by fusion (ACDF).
  \item \textbf{Nerve Root Identification and Retraction:} Under microscopic or endoscopic visualization, the dura mater and the compressed nerve root are carefully identified. The nerve root is gently retracted medially to expose the herniated disc material.
  \item \textbf{Disc Fragment Removal:} The herniated nucleus pulposus fragment, which is typically extruded or sequestered, is carefully grasped with specialized pituitary rongeurs and removed. Any free fragments within the spinal canal are also extracted. The disc space may be explored to ensure complete decompression and remove any loose fragments.
  \item \textbf{Decompression Confirmation:} The nerve root is visually inspected to confirm it is free of compression, mobile, and pulsates normally, indicating successful decompression.
  \item \textbf{Hemostasis and Closure:} Meticulous hemostasis is achieved using bipolar cautery. The muscles are allowed to fall back into place, and the fascia, subcutaneous tissue, and skin are closed in layers. A small drain may be placed temporarily, though this is less common with modern minimally invasive techniques.
\end{itemize}
Minimally invasive techniques, such as microdiscectomy using an operating microscope or endoscopic discectomy, utilize smaller incisions and specialized retractors or endoscopes to minimize muscle trauma, reduce blood loss, and facilitate faster recovery.

\section*{Complications \& Risks}
As with any surgical procedure, discectomy carries potential risks, which can be broadly categorized into general surgical risks and procedure-specific risks.

\textbf{General Surgical Risks:}
\begin{itemize}
  \item \textbf{Bleeding:} Intraoperative blood loss, or postoperative epidural hematoma formation, which can compress neural structures and potentially require reoperation.
  \item \textbf{Infection:} Surgical site infection (superficial or deep), discitis (infection of the disc space), or meningitis (a rare but serious infection of the meninges).
  \item \textbf{Anesthesia Complications:} Adverse reactions to anesthetic agents, respiratory depression, cardiac events (e.g., arrhythmia, myocardial infarction), stroke, or nerve injury from prolonged or improper positioning.
\end{itemize}

\textbf{Procedure-Specific Risks:}
\begin{itemize}
  \item \textbf{Dural Tear with Cerebrospinal Fluid (CSF) Leak:} Accidental puncture of the dura mater, the membrane surrounding the spinal cord and nerve roots, leading to CSF leakage. This may require primary repair and can result in postoperative headaches, pseudomeningocele formation, or, rarely, meningitis.
  \item \textbf{Nerve Root Injury:} Direct trauma to the nerve root during retraction or disc removal, potentially leading to new or worsened weakness, numbness, or pain in the distribution of that nerve.
  \item \textbf{Recurrent Disc Herniation:} The remaining disc material can re-herniate at the same level, occurring in 5-15\% of patients, sometimes requiring repeat surgery.
  \item \textbf{Spinal Instability:} While rare with discectomy alone, excessive bone removal or pre-existing instability can be exacerbated, potentially necessitating a subsequent fusion procedure.
  \item \textbf{Vascular Injury:} Extremely rare but serious, particularly with anterior approaches (e.g., great vessels in the abdomen for lumbar, carotid artery or vertebral artery for cervical).
  \item \textbf{Bowel or Bladder Dysfunction:} Can occur with cauda equina syndrome or direct injury to sacral nerve roots, though rare with standard discectomy.
  \item \textbf{Persistent Pain or Failed Back Surgery Syndrome:} Despite technically successful decompression, some patients may experience ongoing back pain or radicular symptoms, often due to chronic nerve damage, scar tissue formation (epidural fibrosis), or other underlying spinal pathology.
  \item \textbf{Scar Tissue Formation (Epidural Fibrosis):} Scar tissue can form around the nerve root, potentially leading to recurrent symptoms, though its clinical significance and correlation with pain are debated.
  \item \textbf{Wrong Level Surgery:} A rare but devastating complication where surgery is performed at an incorrect spinal level, emphasizing the importance of intraoperative imaging verification.
\end{itemize}

\section*{Postoperative Care \& Recovery}
Immediately after a discectomy, the patient is transferred to the Post-Anesthesia Care Unit (PACU) for close monitoring. Vital signs, neurological status (motor strength, sensation), and pain levels are continuously assessed. Pain management is initiated, typically with intravenous or oral analgesics, aiming for effective control to facilitate early mobilization. Most patients are encouraged to mobilize early, often walking with assistance within hours of surgery, which helps prevent complications like deep vein thrombosis and promotes overall recovery.

Most patients are discharged home the same day or the following day, depending on their recovery progress, pain control, and the type of discectomy performed. Discharge instructions include detailed wound care, a pain medication schedule, and strict activity restrictions. For the first 1-2 weeks, patients are advised to avoid bending, lifting (typically no more than 5-10 pounds), and twisting motions of the spine. Gradual increase in ambulation is encouraged, but prolonged sitting or standing should be limited. Physical therapy may be initiated to teach proper body mechanics, strengthen core muscles, and improve flexibility. Long-term recovery involves a progressive return to activities, with light activities typically resumed within 4-6 weeks and more strenuous activities or heavy lifting restricted for 3-6 months, depending on the individual's healing, the extent of the surgery, and the surgeon's recommendations.

During a discectomy, the surgeon meticulously documents the intraoperative findings, which typically include the precise location and nature of the disc herniation (e.g., protrusion, extrusion, sequestration), the degree of nerve root compression, and any associated inflammatory changes or adhesions. The appearance of the nerve root (e.g., stretched, edematous, discolored) and its pulsatility after decompression are key observations. The size and consistency of the excised disc fragment are also noted.

The resected disc material is sent to pathology for histological examination. The pathology report typically confirms the presence of fragments of nucleus pulposus and annulus fibrosus, consistent with disc herniation. Microscopic examination may reveal degenerative changes within the disc tissue, such as chondroid metaplasia, fibrosis, and sometimes inflammatory cells. The primary purpose of the pathological examination is to confirm the benign nature of the excised tissue and rule out other pathologies, such as infection or malignancy, which are rare but important considerations in differential diagnosis.

A normal and successful postoperative course following discectomy is characterized by a predictable progression of recovery and symptom resolution. Patients typically experience immediate and significant relief of radicular pain (leg pain for lumbar, arm pain for cervical) shortly after surgery, indicating successful decompression of the nerve root. While some incisional pain is expected, it is usually manageable with oral analgesics and gradually subsides over days to weeks. Numbness and motor weakness, if present preoperatively, may take longer to resolve, often improving gradually over several weeks to months as the nerve heals. Patients are encouraged to mobilize early, and wound healing typically progresses without complications. The goal is a progressive return to normal daily activities, with improved mobility and reduced reliance on pain medication, leading to an overall enhancement in functional capacity and quality of life.

Postoperative care and follow-up are essential to monitor recovery, address any concerns, and guide rehabilitation after discectomy.

\begin{itemize}
  \item \textbf{Postoperative Clinic Visits:} Typically scheduled at 1-2 weeks for a wound check and initial assessment of pain and neurological status, followed by visits at 6 weeks, 3 months, and sometimes 6 months or a year, depending on the patient's progress and the surgeon's discretion.
  \item \textbf{Suture/Staple Removal:} If non-absorbable sutures or staples were used for skin closure, they are usually removed at the first follow-up visit.
  \item \textbf{Physical Therapy (PT) Referral:} Many patients are referred for a structured physical therapy program, typically starting a few weeks post-op, to strengthen core muscles, improve flexibility, and learn proper body mechanics to prevent future injury and optimize long-term spinal health.
  \item \textbf{Activity Resumption Guidance:} The surgeon and physical therapist will provide detailed guidance on gradually increasing activity levels, including return to work, exercise, and sports, tailored to the individual's recovery.
  \item \textbf{Imaging Studies:} Routine postoperative imaging (MRI, X-rays) is generally not required unless there are persistent or recurrent symptoms, new neurological deficits, or concerns for complications like infection or re-herniation.
\item \textbf{Medication Management:} Gradual tapering of pain medications and anti-inflammatories as symptoms improve.
\end{itemize}
Patients are educated on warning signs to call the doctor immediately, such as fever, worsening radicular pain, new or progressive weakness or numbness, bowel/bladder changes, or excessive wound drainage or redness.

\section*{Outcomes \& Prognosis}
Discectomy is one of the most frequently performed spinal surgical procedures worldwide, primarily addressing symptomatic intervertebral disc herniation. Lumbar disc herniation, the most common indication for discectomy, has an estimated annual incidence of 5 to 20 cases per 1,000 adults, with a peak incidence occurring between the ages of 30 and 50 years. There is a slight male predominance. Cervical disc herniation is less common but also a significant cause of radiculopathy and myelopathy, with an incidence of approximately 1 per 10,000 adults per year. The overall mortality rate associated with discectomy is extremely low, typically less than 0.1\%. Major morbidity rates are also low, with complications such as dural tear occurring in 1-5\% of cases, and surgical site infection in less than 1\%. The high prevalence of disc-related pain and the effectiveness of discectomy in carefully selected patients contribute to its widespread application.

The outcomes following discectomy are generally excellent, particularly for the relief of radicular pain. Success rates for significant improvement in leg or arm pain range from 80\% to 95\%. Functional outcomes, including return to work and daily activities, are also highly favorable for the majority of patients. While radicular pain often resolves quickly, improvement in numbness and weakness may be slower and less complete, especially if symptoms were long-standing or severe preoperatively. The landmark Spine Patient Outcomes Research Trial (SPORT) demonstrated that surgical treatment for lumbar disc herniation resulted in greater improvement in pain and function compared to non-operative treatment at 2-year follow-up, particularly for patients with clear indications.

The recurrence rate for disc herniation at the same level after discectomy is approximately 5-15\% over 5-10 years, with some studies suggesting higher rates in specific patient populations (e.g., smokers, those with large annular defects). Prognostic factors influencing outcomes include the duration of preoperative symptoms (shorter duration often correlates with better outcomes), the severity of neurological deficit, the presence of inflammatory markers, and patient-specific factors like smoking status, obesity, and psychological well-being. Although most patients experience lasting relief, a subset may develop persistent axial back or neck pain, or require further surgical intervention for recurrent herniation or adjacent segment disease over time.

\section*{References}
\begin{enumerate}
  \item MedlinePlus. Diskectomy. U.S. National Library of Medicine; 2025.
  \item Chapman JR, et al. Spinal Instrumentation. In: Campbell's Operative Orthopaedics. 14th ed. Elsevier; 2021.
  \item North American Spine Society (NASS). Evidence-Based Clinical Guidelines for Multidisciplinary Spine Care. Lumbar Discectomy. 2020.
  \item Greenberg MS. Handbook of Neurosurgery. 9th ed. Thieme; 2020.
\end{enumerate}

\clearpage
